% Intended pattern for every results table: caption and \label live here, the body is
% generated by the evaluation scripts and pulled in with \input. Reference it with \cref.
%
%   \begin{table}
%     \centering
%     \caption{Wall precision, recall and F1 on held-out drawings.}
%     \label{tab:baseline}
%     \input{../results/tables/baseline.tex}
%   \end{table}
%
%   ... as \cref{tab:baseline} shows ...
\chapter{Results}
\label{ch:results}

\section{Data Preparation Findings}
\label{sec:data-preparation-findings}
% source: runs/2026-10-26-floorplancad-parse/metrics.json
The parsing run establishes a public inventory of distinct drawing contents and
their entities, with the coverage and type counts reported in
\cref{tab:public-parse-summary,tab:public-entity-kinds}.
% source: runs/2026-10-26-floorplancad-parse/metrics.json
No file failures are reported in that run.
% source: docs/decisions/0010-evidence-backed-thesis-draft.md
This finding concerns data preparation and does not measure whether the semantic
classes can be learned from geometry.

% source: runs/2026-10-26-floorplancad-split-audit-2/metrics.json
The source audit identifies exact-duplicate content clusters spanning the official
training and test folders, together with a large share of official test files
whose source group also occurs in training, as shown in
\cref{tab:public-source-audit}.
% source: docs/decisions/0006-floorplancad-split-by-source-group.md
The official release folders therefore do not provide the source separation
required by the adopted evaluation design.
% source: docs/decisions/0006-floorplancad-split-by-source-group.md
The experiments use deduplication and source-group assignment instead.

\begin{table}[htbp]
  \centering
  \small
  % source: runs/2026-10-26-floorplancad-split-audit-2/metrics.json
  \caption{Source overlap and exact-duplicate evidence in the official release
  folders, before the source-group split used in this work.}
  \label{tab:public-source-audit}
  \input{tables/public_source_audit.tex}
\end{table}

% source: runs/2026-10-26-floorplancad-split-audit-2/metrics.json
Layer-vocabulary similarity also distinguishes the release folders that share a
numbering namespace from the separately numbered folder, as reported in
\cref{tab:source-layer-similarity}.
% source: docs/decisions/0006-floorplancad-split-by-source-group.md
The same-prefix comparison supports the adopted namespace grouping, while the
random-pair comparison provides a reference for interpreting those similarities.
% source: docs/decisions/0006-floorplancad-split-by-source-group.md
A prefix alone is insufficient to identify a common source across all folders.

\begin{table}[htbp]
  \centering
  \small
  % source: runs/2026-10-26-floorplancad-split-audit-2/metrics.json
  \caption{Layer-vocabulary Jaccard similarities recorded by the source audit.
  Similarity supports a grouping convention, but does not verify building identity.}
  \label{tab:source-layer-similarity}
  \input{tables/public_layer_similarity.tex}
\end{table}

% source: docs/decisions/0007-floorplancad-fold-assignment.md
The resulting training, validation and held-out assignments are fixed in the split
run, so later experiments can evaluate on the same drawing partition.
% source: runs/2026-10-26-floorplancad-split/metrics.json
The recorded folds contain each target class, with walls more numerous than doors
or windows in every fold, as shown in \cref{tab:public-class-counts}.
% source: docs/decisions/0006-floorplancad-split-by-source-group.md
These checks do not exclude near-duplicate crops filed under different groups,
which remains a limit of the current source separation.

\section{Baseline}
\label{sec:baseline}
\todo[inline]{No trained baseline result is recorded. Introduce the baseline with a method
decision, then include its generated wall precision, recall and F-score table from
the corresponding training and evaluation runs. Do not substitute dataset-audit
results for model performance.}

\section{Comparison of Representations}
\label{sec:representation-comparison}
\todo[inline]{Include the generated representation comparison after both branches have
been evaluated on the same held-out entities. Report the raster-to-vector mapping
and geometry-only inputs with the comparison; report doors and windows separately
from walls.}

\section{Public Data versus the Archive}
\label{sec:public-vs-archive}
\todo[inline]{Compare public-data evaluation with external archive evaluation only after
conversion and reference-label verification have been recorded. The current design
trains on public data only. No archive training or fine-tuning result is planned
under the current decision.}

\section{Failure Analysis}
\label{sec:failure-analysis}
\todo[inline]{Include generated failure counts and approved aggregate breakdowns tied to
evaluation run identifiers. Distinguish classification errors from parsing,
conversion, label-recovery and wall-generation failures where the recorded outputs
support that distinction.}
